\documentclass{article}
\usepackage[utf8]{inputenc}
\usepackage[spanish]{babel}
\usepackage{listings}
\usepackage{graphicx}
\graphicspath{ {images/} }
\usepackage{cite}

\begin{document}

\begin{titlepage}
    \begin{center}
        \vspace*{1cm}
            
        \Huge
        \textbf{Taller Memoria}
            
        \vspace{0.5cm}
        \LARGE
        Trabajo de investigación.
            
        \vspace{1.5cm}
            
        \textbf{Juan Camilo Mazo Castro}
            
        \vfill
            
        \vspace{0.8cm}
            
        \Large
        Despartamento de Ingeniería Electrónica y Telecomunicaciones\\
        Universidad de Antioquia\\
        Medellín\\
        Septiembre de 2020
            
    \end{center}
\end{titlepage}

\tableofcontents

\section{Introducción}
En el presente trabajo se hace con el fin de entender mejor el funcionamiento y la gestión que se le da a la memoria de un computador, esta información será fundamental para poder entender algunos conceptos y conocimientos necesarios que debe tener un programador. Se resolverán una serie de preguntas planteadas con respecto al tema comentado anteriormente.

\section{Sección de contenido} \label{contenido}

1.Defina qué es la memoria de un computador.

Una memoria es un dispositivo electrónico que sirve para almacenar información ya sea temporal o permanentemente. Normalmente se utiliza el termino memoria para referirse a dispositivos de almacenamiento temporal y de alta velocidad de acceso como la memoria RAM o las memorias Cache incluidas en algunos procesadores, sin embargo, también hay otras de almacenamiento permanente como los discos duros.
Las memorias cumplen un papel muy importante en el funcionamiento de las maquinas. Éstas junto con otros dispositivos permiten que puedan funcionar correctamente todos los programas que utilizamos hoy en día, permiten que se almacene y se procese toda la información que solicitamos y requerimos. Las memorias tienen una capacidad de almacenamiento de la información mucho más grande que otros dispositivos, por lo que se hacen indispensables para que nuestras computadoras funcionen. Para lograr el objetivo antes dicho, los diferentes tipos de memoria trabajan en equipo junto con otros dispositivos como el microprocesador y los controladores.

2.Mencione los tipos de memoria que conoce y haga una pequeña descripción de cada tipo.

Disco Duro: Tiene una gran capacidad de almacenamiento y funciona con unos discos que giran, es más lenta para dar acceso a la información por lo que es utilizada para almacenar datos que se quieren guardar permanentemente. Es muy importante pues guarda información que se requiere para que el usuario pueda hacer uso de la computadora, los datos necesarios para efectuar cualquier programa se toman de esta y se ponen en la memoria RAM para que el microprocesador tenga un acceso más rápido a ellos. Además, el Disco Duro guardará permanentemente todas las modificaciones que se le hagan a los archivos que tiene el usuario cuando el desee guardarlos.

Memoria Virtual: Es una parte del Disco Duro que se utiliza para guardar datos que ocupan una buena parte de almacenamiento y que no se utilizan tanto, pero que de igual forma se mantienen ahí para que se puedan utilizar cuando sean requeridos y mientras se encuentran en la Memoria Virtual no ocupan espacio innecesario en la memoria RAM.

Memoria RAM: es una memoria de almacenamiento temporal que es de gran ayuda para el microprocesador, mientras se ejecuta un programa esta almacena los datos temporales que el microprocesador va requiriendo a medida que ejecuta las ordenes que se le demandan. Tiene menor capacidad de almacenamiento y da acceso a la información más rápido que un Disco Duro. La memoria RAM almacena también las ordenes que se le solicitan al microprocesador para que este las pueda leer y ejecutar.
 
Memoria Cache: Es una memoria de menor capacidad que la memoria RAM, se encuentra en algunos procesadores como los Intel Core i5 y al encontrarse en el mismo procesador, tiene un acceso más rápido a la información que estas guardan. Existen tres tipos de Caché diferentes, la L1, L2 y L3. L1 es la más veloz y es la que tiene menor almacenamiento de las tres,L2 es un punto intermedio entre L1 Y L3, L3 tiene el mayor almacenamiento de las tres pero es menos veloz que L1 Y L2, aun así sigue siendo más veloz que la memoria RAM. Se utiliza para almacenar datos que se utilizan con mayor frecuencia por el usuario, así le evita al procesador tener que retirar la información de la memoria RAM que es más lenta.

Memoria USB: Es una memoria portable de almacenamiento permanente, no tienen demasiado espacio de almacenamiento, pero sirven para guardar datos de un computador y utilizarlos en otro.

Memoria ROM: Es una memoria no volátil que se encuentra en la placa madre de los computadores, esta le da las primeras instrucciones al microprocesador cuando la maquina se enciende.

DRAM: Es un tipo de memoria volátil, es decir que cuando se le corta la corriente eléctrica los datos en ella desaparecen. Se le llama Dynamic RAM porque cada cierto tiempo necesita un ciclo de refresco de datos para que estos permanezcan, incluso cuando la maquina se encuentra encendida.\cite{uruguayoc}


3.Describa la manera en que se gestiona la memoria en el computador.

Cuando nuestras computadoras se encienden lo primero que hace el procesador es esperar que se le indiquen las primeras instrucciones, estas las tiene la memoria ROM que guarda de manera permanente las instrucciones que se deben realizar siempre que empieza a trabajar la máquina. Una vez la memoria RAM ya se encuentra en funcionamiento lo primero que hace es cargar la BIOS que contiene información del sistema operativo que se encuentra en el Disco Duro de arranque, así una vez iniciado el sistema operativo el usuario puede empezar a utilizar la computadora. Ya en este punto el usuario empieza a efectuar ciertas ordenes que le dan instrucciones al procesador, el procesador utiliza la memoria RAM para guardar y tener un acceso más rápido a los datos que el usuario requiere para que funcionen los programas que él necesita. Los controladores de memoria hacen un espacio en la memoria RAM para traer del Disco Duro los datos requeridos, una vez en la RAM el procesador hace uso de ellos y cuando ya no los necesita, si tienen modificaciones y el usuario desea guardarlos, el Disco Duro se encarga de guardarlos permanentemente. La memoria RAM siempre carga todos los datos necesarios para que los programas que utilizamos puedan funcionar, cuando los datos son demasiado grandes y la memoria RAM no puede con todos ellos se hace uso de la memoria virtual, pero cuando se hace esto el computador empieza a funcionar de una forma más lenta, pues el procesador tarda más en tener acceso a los datos que se encuentran ahí. Además de la memoria RAM también se cuenta con la memoria Cache que se encuentra en el procesador, esta almacena información que se usa con más recurrencia para que cada vez que el procesador requiera esos datos para ejecutar un programa tenga un acceso más rápido a ellos. En conclusión, la memoria RAM sirve como apoyo para el microprocesador, le ayuda a cargar toda la información que este no puede almacenar por si solo, información que este requiere para efectuar los procesos de las ordenes que se le indican.  La gestión que se le da a cada una de las memorias se hace con el objetivo de tener una mayor eficiencia en la ejecución de los diferentes trabajos que hace la máquina.

4.¿Qué hace que una memoria sea más rápida que otra?¿Por qué es tan importante?

Una memoria se hace más rápida cuando tiene menos capacidad de almacenamiento, es decir que guarda menos información. Esto se hace importante para que nuestras computadoras funcionen más rápido, con una mayor eficiencia y con mayor fluidez. Se requiere que el microprocesador tenga un acceso rápido a los datos que necesita para ejecutar los programas. Por lo antes dicho es necesario emplear dispositivos como la memoria RAM que tiene una capacidad de almacenamiento menor que el Disco Duro, pero que puede guardar temporalmente los datos que el procesador requiere mientras utilizamos nuestro computador. También se utilizan por ejemplo las memorias Cache que se encuentran en el mismo procesador y funcionan para almacenar los datos más recurrentes en el uso diario que le da el usuario a la máquina. Todo este trabajo en equipo facilita las tareas que debe hacer el procesador y permite que se pueda hacer un trabajo más rápido.

\begin{lstlisting}
#include <stdio.h>
#define N 10
/* Block
 * comment */

int main()
{
    int i;

    // Line comment.
    puts("Hello world!");
    
    for (i = 0; i < N; i++)
    {
        puts("LaTeX is also great for programmers!");
    }

    return 0;
}
\end{lstlisting}

A continuación se presenta el logo de C++ Figura (\ref{fig:cpplogo})

\begin{figure}[h]
\includegraphics[width=4cm]{cpplogo.png}
\centering
\caption{Logo de C++}
\label{fig:cpplogo}
\end{figure}

En la sección de teoremas (\ref{contenido})

\section{Conclusión} \label{conclulsion}

\bibliographystyle{IEEEtran}
\bibliography{references}

\end{document}
